% Galería de las plantillas de diagrama (assets/*.tex), una por slide. La usa el README.
\documentclass[aspectratio=169,11pt]{beamer}
\def\tema{bloques}\def\idioma{spanish}\def\nbloques{1}\def\logos{}
\def\estilodir{../../../charlas-charlie/assets/estilo/}
\makeatletter\def\input@path{{../../../charlas-charlie/assets/estilo/}{../../assets/plantillas/}}\makeatother
\input{preambulo}
\input{selector.tex}
\input{chips.tex}
\begin{document}

\begin{frame}{flujo.tex · un foco, borde negro, bucle}\centering\input{flujo.tex}\end{frame}
\begin{frame}{pipeline.tex · cadena y salidas}\centering\input{pipeline.tex}\end{frame}
\begin{frame}{selector.tex · una opción por columna}\centering\selector{imp}{oss}{din}{vreman}{imp oss din + vreman}\end{frame}
\begin{frame}{chips.tex · tabla de casos}\centering
  \tablacasos{%
    \caso{imp asgs din}\nota{4 mallas} & & \chip{serieB}{ASGS} & \chip{serieC}{din} & \\
    \caso{imp oss din}                 & & \chip{serieB}{OSS}  & \chip{serieC}{din} & \\
    \caso{exp}                         & \chip{serieA}{EXP} &               &                    & \\
    \caso{exp + vreman}                & \chip{serieA}{EXP} &               &                    & \chip{serieD}{VREMAN}\\}
\end{frame}
\begin{frame}{imagen\_anotada.tex · zona resaltada y flecha}\centering\input{imagen_anotada.tex}\end{frame}
\begin{frame}{ejes.tex · idea cualitativa con ejes}\centering\input{ejes.tex}\end{frame}
\begin{frame}{matriz.tex · matriz por bloques}\centering\input{matriz.tex}\end{frame}
\begin{frame}{llamadas.tex · secuencia de llamadas}\centering\input{llamadas.tex}\end{frame}
\begin{frame}{esquema\_fisico.tex · geometría con cotas}\centering\input{esquema_fisico.tex}\end{frame}
\begin{frame}{elemento.tex · elemento de malla}\centering\input{elemento.tex}\end{frame}
\end{document}
